Nous avons appris au chapitre \ref{ch:proc-pascal} à déclarer et à appeler des procédures sans paramètres
nous permettant de répéter plusieurs fois la même séquence de commmandes du
\robot\ à des endroits différents dans le programme. Au chapitre \ref{ch:appel} nous avons
utilisé des commandes du \robot\ qui sont des appels de procédures avec paramètres,
l'usage d'un ou plusieurs paramètres per\-met\-tant à la procédure de faire des séquences
de commandes similaires mais différentes selon la valeur des paramètres. Nous allons maintenant
voir comment déclarer des procédures avec paramètres en \ocaml. Commençons
par un exemple simple~:

\bigskip

 {\raggedright 
 {}\hfill\vbox to 0pt{\vspace*{\baselineskip}\vspace*{3cm}
 \includevspacePostScript[-\baselineskip]{CarrePointille}{0.3}{0.3}}\hspace*{30mm}%
 \label{prog:CarrePointille}
 \input ocaml/CarrePointille.tex
 }
La procédure {\idt traitPointille} permet de tracer un trait en pointillé cons\-ti\-tué par une succession de
{\idt n} pixels noirs séparés par $(${\idt n}$ - 1)$ pixels blancs, le nom\-bre {\idt n} de points étant
passé en paramètre de cette procédure. Dans la déclaration `{\mcl let}\ {\idt traitPointille\ n\ =}',
le paramètre {\idt n} est écrit après le nom de la procédure, à la place des parenthèses
{\idt ()} des procédures sans paramètres. Il n'est pas nécessaire d'indiquer que {\idt n} est
un nombre entier~: le compilateur \ocaml\ le découvre tout seul, en voyant que {\idt n}
est utilisé dans l'expression entière {\idt n - 1}.
La procédure {\idt carre} permet de dessiner un carré
en pointillé, le nombre {\idt p} de points par côté étant passé en paramètre. L'appel {\idt carre\ 30}
indique qu'il faut exécuter la procédure {\idt carre} avec un paramètre {\idt p} égal à 30. De ce fait la
procédure {\idt traitPointille} est appelée quatre fois avec un paramètre {\idt n} égal à {\idt p} soit
30. On dit que {\idt p} est le {\em paramètre formel}\index{Paramètre formel}  (dont on ne connaît pas encore la valeur
au moment de la déclaration de la procédure et que l'on désigne par un identificateur) et que 
30 est le {\em paramètre effectif}\index{Paramètre effectif} (c'est-à-dire
la valeur du paramètre formel pour cet appel de la procédure).

De manière générale, les déclarations et appels de procédures avec pa\-ra\-mè\-tres ont  la forme
suivante~:

\index{Paramètre formel}\index{Déclaration de procédure}\index{Identificateur}
\begin{center}
\begin{tabular}{@{}|p{12.5cm}|@{}} \hline
{\em Déclaration de procédure avec paramètres~:}\\ \hline \hline
\begin{minipage}[t]{12.5cm}
{\medskip
\mbox{\blc\ \ \ \cmt (*\ \small\rm Spécification : ce que fait la procédure\ \cmt *)}\par 
\mbox{\blc\ \ \ \mcl let\blc\hspace{0.3em}\blanc\hspace{0.3em}\idt nomDeLaProcedure%
\ paramFormel1\footnote{\ Un paramètre formel est un identificateur.}\ paramFormel2\ \dots\ =}\par 
\mbox{\blc\ \ \ \ \ \small\rm \ldots\ Instructions du corps de la procédure%
\footnote{\ Les expressions  figurant dans ces instructions peuvent utiliser les paramètres formels
{\idt paramFormel1}, \ldots\ de la procédure.}
 (séparées par des `$;$') \ldots}\par 
\medskip}
\end{minipage}\\ \hline 
\end{tabular}
\end{center}

\index{Paramètre effectif}\index{Appel de procédure}
\begin{center}
\begin{tabular}{@{}|p{12.5cm}|@{}} \hline
{\em Appel de procédure avec paramètres~:}\\ \hline \hline
\begin{minipage}[t]{12.5cm}
{\medskip
\mbox{\blc\ \ \ \idt nomDeLaProcedure\ paramEffectif1%
\footnote{\ Un paramètre effectif est une expression entière (ou une couleur, etc.), 
placée entre parenthèses si ce n'est pas un simple nombre ou identificateur.}\ paramEffectif2\ \dots%
\footnote{\ Les paramètres effectifs sont donnés dans l'ordre où apparaissent les pa\-ra\-mè\-tres
formels correspondants.}}\par 
\medskip}
\end{minipage}\\ \hline 
\end{tabular}
\end{center}

On remarquera que les paramètres, formels dans la déclaration de procédure comme effectifs
dans l'appel de procédure, sont simplement séparés par des blancs, sans parenthèses ni
virgules. Par exemple pour dessiner un rectangle de hauteur 3 et de largeur 4 avec la 
procédure {\idt rectangle} du programme dans le dossier 
{\idt livre/\allowbreak 13\_decl\_procedure\_avec\_parame}, on écrit {\idt rectangle\ 3\ 4}.
Par contre un paramètre effectif qui est une expression doit être placé entre parenthèses,
comme dans {\idt demiCercleDroit\ (2\ *\ tailleGrille)} (page \pageref{prog:HeliceEnCouleur})
ou {\idt avf\ (-3)}. Sans les parenthèses, {\idt demiCercleDroit\ 2\ *\ tailleGrille} serait
compris comme {\idt (demiCercleDroit\ 2)\ *\ tailleGrille} c'est-à-dire comme la 
multiplication par {\idt tailleGrille} du résultat de l'appel {\idt demiCercleDroit\ 2}, 
ce qui n'a pas de sens. Le compilateur le signale par le message~:
`{\tt This expression has type unit but an expression was expected of type int}'
(cette expression est de type {\idt unit}, celui des instructions, alors qu'une 
expression de type entier était attendue).


Une procédure avec 
paramètres formels peut être appelée plusieurs fois dans un programme avec des 
pa\-ra\-mè\-tres effectifs
qui peuvent être différents à chaque fois.

Les paramètres formels d'une procédure ne sont utilisables que dans le
corps de la procédure. Par conséquent des procédures différentes peuvent avoir des
paramètres formels portant le même nom bien qu'il s'agisse de paramètres différents. 

Ces différents 
points sont illustrés dans l'exemple ci-dessous.

\begin{exemple}[Spirale segmentée]{\rm
~Le programme ci-dessous dessine une spirale constituée de paires de segments de droite
de longueurs successives 1, 2, 3, \ldots, 19~:

\newpage
 {\raggedright 
 \mbox{}\hspace*{40mm}\vbox to 0pt{\vspace*{\baselineskip}\vspace*{4cm}
\includevspacePostScript[-\baselineskip]{SpiraleSegmentee}{0.7}{0.7}}
 \input ocaml/SpiraleSegmentee.tex
 }
\noindent $\Box$ }
\end{exemple}


\noindent \rule{137mm}{0.1mm} %---------------------------------------------
\begin{exercices}{Écrire des programmes \ocaml\ comportant des procédures avec
paramètres pour réaliser les dessins suivants~:

\tiret la croix à trois branches transversales, l'hélice coloriée, le mi\-cro-\-or\-di\-na\-teur,
le logo en couleur, la chenille (où le nombre
d'anneaux du corps de la chenille est passé en paramètre de la procédure de dessin), la pirogue 
(où le nombre de rameurs est un paramètre) représentés à la page \pageref{fig:crx3b}$\:$;

\tiret les carrés inscrits les uns dans les autres
ou la version peinte de ce dessin représenté à la figure \ref{fig:carinspt} (page \pageref{fig:carinspt})$\:$;

\tiret l'arabesque inspirée par la coupole du mausolée du Sultan
         Qayt-\-Bay re\-pré\-sen\-tée à la figure \ref{fig:arab2} (page \pageref{fig:arab2})$\:$;

\tiret la péniche de la figure \ref{fig:transports} (page \pageref{fig:transports}) où le nombre
         de hublots est donné en paramètre$\:$;

\tiret le pavage de Golomb de la figure \ref{fig:golomb} (page \pageref{fig:golomb}).

\noindent $\Box$}\end{exercices}
\twocolumn
\begin{figure}[h]
%\hspace{-3.363cm} 
\includevspacePostScript[4.9cm]{Arabesque2}{0.95}{0.95}
 \caption{Arabesque}\label{fig:arab2}
\end{figure}
\vfill
\begin{figure}[h]
%\hspace{-3.363cm} 
\includevspacePostScript[0.0cm]{HeliceEnCouleur}{0.95}{0.95}
 \caption{Hélice coloriée}
 %\label{fig:helice}
\end{figure}
\vfill
\begin{figure}[h]
%\hspace{-3.363cm} 
\includevspacePostScript[0.0cm]{MacintoshSE}{0.95}{0.95}
 \caption{Ancien micro-ordinateur}
 %\label{fig:MacSE}
\end{figure}
\mbox{}
\newpage
\begin{figure}[!h]
%\hspace{-3.363cm} 
\includevspacePostScript[5.3cm]{LogoEnCouleur}{0.95}{0.95}%
\caption{Logo en couleur}%
%\label{fig:Logoc}
\end{figure}
\begin{figure}[!h]
%\hspace{-3.363cm} 
\includevspacePostScript[1.0cm]{Chenille}{0.95}{0.95}%
\caption{Chenille}%
%\label{fig:chenille}
\end{figure}
\begin{figure}[!h]
%\hspace{-3.363cm} 
\includevspacePostScript[1.0cm]{Pirogue_1}{0.95}{0.95}%
\caption{Pirogue}%
%\label{fig:pirog1}
\end{figure}
\vfil
\begin{figure}[!h]
\includevspacePostScript[0.0cm]{Croix3branches}{0.95}{0.95}%
\caption{Croix à trois branches}%
\label{fig:crx3b}%
\end{figure}
\onecolumn
\begin{figure}[h]
\includevspacePostScript[5.7cm]{CarresInscritsPeints}{0.95}{0.95}
 \caption{Motif de carrés inscrits de Horemis}\label{fig:carinspt}
\end{figure}
\begin{corrige}{
Nous donnons un programme, les autres se trouvant dans le dossier 
{\idt livre/\allowbreak 13\_decl\_procedure\_avec\_parame}~:
\label{prog:HeliceEnCouleur}
\input ocaml/HeliceEnCouleur.tex
\noindent $\Box$}\end{corrige}

\vfill

%%%%%%%%%%%%%%%%%%%%%%%%%%%%%%%%%%%%%%%%%%%%%%%%%%%%%%%%%%%%%%%%%%%%%%%%%%%
\noindent \framebox[\textwidth]{
\begin{minipage}{\myparboxwidth}
\setlength{\parindent}{2em}
\small \sf
\vspace{\abovedisplayskip}
\section*{Addenda}

On distingue en \ocaml, les définitions des identificateurs de leurs utilisations. \index{Identificateur}
Les   {\em occurrences de définition} figurent dans une déclaration {\mcl let} de
constante ou de procédure comme nom de constante, de procédure ou de paramètre formel.
Les  {\em occurrences d'utilisation} de ces identificateurs sont les noms de constante ou de
paramètre formel utilisés dans une expression entière ou les noms de procédure utilisés dans des appels.

 Il se peut qu'un même identificateur ait plusieurs occurrences de définition dans un programme. De ce fait il se
 pose le problème d'établir la correspondance entre les occurrences d'utilisation et de définition d'un même
 identificateur. Par 
exemple l'identificateur peut être utilisé pour nommer une constante du programme et des paramètres
formels dans des procédures. En \ocaml, ces procédures peuvent être imbriquées les unes dans les
 autres (en déclarant une procédure {\idt p2} à l'intérieur du corps d'une procédure {\idt p1}
 avec {\mcl let}\ \ldots\ {\mcl in}), comme dans l'exemple ci-après~:
 
\smallskip

{
\parindent \parindentprog
\medskip

\vspace{\neglskip}\mbox{\blc {\mcl open}{\idt \ LibRobot}}\par
\vspace{0.3\baselineskip}
\vspace{\neglskip}\mbox{\blc {\mcl let}{\idt \ n\ =\ 17}}\par
\vspace{0.3\baselineskip}
\vspace{\neglskip}\mbox{\blc {\mcl let}{\idt \ p1\ n\ =}}\par
\vspace{\neglskip}\mbox{\blc \blk{\mcl let}{\idt \ p2\ n\ =}}\par
\vspace{\neglskip}\mbox{\blc \blk\blk{\idt ...}}\par
\vspace{\neglskip}\mbox{\blc \blk\blk{\cmt (*1*)}{\idt \ ...\ n\ ...}}\par
\vspace{\neglskip}\mbox{\blc \blk\blk{\idt ...}}\par
\vspace{\neglskip}\mbox{\blc \blk{\mcl in}}\par
\vspace{\neglskip}\mbox{\blc \blk{\idt ...}}\par
\vspace{\neglskip}\mbox{\blc \blk{\cmt (*2*)}{\idt \ p2\ n;}}\par
\vspace{\neglskip}\mbox{\blc \blk{\idt ...}}\par
\vspace{0.3\baselineskip}
\vspace{\neglskip}\mbox{\blc {\mcl let}{\idt \ ()\ =}}\par
\vspace{\neglskip}\mbox{\blc \blk{\idt ...}}\par
\vspace{\neglskip}\mbox{\blc \blk{\cmt (*3*)}{\idt \ p1\ n;}}\par
\vspace{\neglskip}\mbox{\blc \blk{\idt ...}}\par

\medskip }

On dit que le programme a une {\em structure de blocs}. \index{Structure de blocs}Ces blocs sont formés par les procédures imbriquées 
 les unes dans les autres, le programme étant  considéré comme le bloc le plus externe. 
 
 La {\em règle de 
 visibilité}\index{Visibilité} des identificateurs en \ocaml\ dit que l'occurrence d'utilisation d'un identificateur
 se réfère à l'occurrence de définition du bloc le plus interne où figure cette occurrence d'utilisation et
 qui contient une occurrence de définition de l'identificateur. 

\vspace{\abovedisplayskip}
\end{minipage}
}
%%%%%%%%%%%%%%%%%%%%%%%%%%%%%%%%%%%%%%%%%%%%%%%%%%%%%%%%%%%%%%%%%%%%%%%%%%%% %%%%%%%%%%%%%%%%%%%%%%%%%%%%%%%%%%%%%%%%%%%%%%%%%%%%%%%%%%%%%%%%%%%%%%%%%%%
%%%%%%%%%%%%%%%%%%%%%%%%%%%%%%%%%%%%%%%%%%%%%%%%%%%%%%%%%%%%%%%%%%%%%%%%%%%
\noindent \framebox[\textwidth]{
\begin{minipage}{\myparboxwidth}
\setlength{\parindent}{2em}
\small \sf
\vspace{\abovedisplayskip}
 
 
Dans l'exemple ci-dessus,
\begin{list}{$\scriptscriptstyle\bullet$}{
}
 \item l'occurrence d'utilisation de {\idt n} à la ligne 1 se réfère au paramètre formel de la procédure {\idt p2}$\:$;
\item l'occurrence d'utilisation de {\idt n} à la ligne 2 se réfère au paramètre formel de la procédure {\idt p1}$\:$;
 \item l'occurrence d'utilisation de {\idt n} à la ligne 3 se réfère à la constante 17.
\end{list}

La structure de blocs, que l'on retrouve dans \ocaml, a été introduite par {\sc algol 60} \cite{algol60}. Ce langage fut créé par un groupe
d'experts européens et américains  réunis dans le groupe de travail WG2.1 de l'IFIP (International Federation of 
Information Processing) qui regroupe toutes les associations nationales d'informaticiens. C'est l'ancêtre d'une famille
de langages parmi lesquels on compte \pascal, {\sc Ada} et {\sc Modula}.

\vspace{\abovedisplayskip}
\end{minipage}
}
%%%%%%%%%%%%%%%%%%%%%%%%%%%%%%%%%%%%%%%%%%%%%%%%%%%%%%%%%%%%%%%%%%%%%%%%%%%%